\scanpage{70}
To describe the action of a $u\in\mathbf{\Gamma}$ (resp $u\in\Gamma$) on $\widehat\Pi_1(X)$ (resp. $\deleted{\unsure{\Pi}}\widehat\Pi_1(X_{\bar K})$), we shall take as ``generators'' of the groupoid $\widehat\Pi_1$ paths coming from $\Pi_1$. We shall have, for ${l}:x\to y$ in $\Pi_1$, ${}^u{l}:{}^ux\to{}^uy$.

If \deleted{\unsure{we assume}} in what follows, for simplicity, we restrict ourselves to rational points $x,y\ldots$ on $X$ -- i.e. we work in the subgroupoid $\widehat\Pi_1(X_K)$ formed by the objects which are invariant under $\mathbf{\Gamma}$ -- then we shall have
\[
(*)\qquad{}^u{l}:{}^ux=x\longrightarrow{}^uy=y,
\]
and we can write, \deleted{f}..
\[
(8)\qquad\begin{cases}
{}^u{l}=g_{l}(u)\circ{l}={l}\circ f_{l}(u)\\
g_{l}(u)\in\widehat\pi_1(X,y)\\
f_{l}(u)\in\widehat\pi_1(X,x)
\end{cases}
\]
\deleted{$g_{l}(u)={}^u{l}\circ{l}^{-1}$}
\figure{\begin{tikzpicture}[x=2mm,y=2mm]
\node at(0,0){,};
\fill(.5,.8)circle(.07);\fill(.75,.55)circle(.07);\fill(1.25,.4)circle(.08);\fill(1.75,.08)circle(.045);
\fill[gray](1.16,-.08)circle(.04);\fill[gray](1.38,-.1)circle(.035);
\end{tikzpicture}}
\margin{We shall also sometimes write $g_u({l})$ instead of $g_{l}(u)$, $f_u({l})$ for $f_{l}(u)$.}
where $f_{l}(u),g_{l}(u),{}^u{l}$ determine one another, \deleted{$f_l(u)$ and $g_l(u)$} by the formulas (8) and
\[
(9)\qquad\begin{cases}
f_{l}(u)={l}^{-1}\cdot{}^u{l},\quad g_{l}(u)={}^u{l}\cdot{l}^{-1}\\
g_{l}(u)={l}(f_{l}(u)),\quad f_{l}(u)={l}^{-1}(g_{l}(u))
\end{cases}
\]
We shall describe the action of a $u\in\Gamma$ \inserted{on a full subgroupoid of $\Pi_1(X)$} by means of the elements $f_{l}(u)$, or $g_{l}(u)$, for a suitable system of generators of the corresponding \inserted{full groupoids} of $\Pi_1(X)$. In the case of a
\margin{69}
